Deux aimants droits identiques sont disposés comme indiqué sur le schéma
ci-dessous. Chaque aimant crée au point $M$ un champ magnétique de valeur
$\qty{2,5e-3}{\tesla}$.
\begin{figure}[H]
    \centering
    \includegraphics[width=0.90\linewidth]{2025-04-20_231812-}
\end{figure}

\begin{enumerate}
    \item Tracer, en précisant l'échelle, les champs $\vec{B}_1$, $\vec{B}_2$ et
          $\vec{B}=\vec{B}_1+\vec{B}_2$.
    \item Déduire de la construction vectorielle la valeur du champ magnétique
          résultant en $M$.
    \item Retrouver le résultat précédent en utilisant une fonction
          trigonométrique.
    \item Comparer la précision des deux résultats.
\end{enumerate}

